
While we have already given some of the following definitions in the
introductory section on the axioms of quantum mechanics, we reproduce
them here for completeness.

\section{Adjoint operators}

If \(A\colon\Cal H\to\Cal H\) is a \emph{bounded} operator,
%then for every
%\(\psi\in \Cal H\), the operator \(f_{\psi}\circ A\cl\Cal H\to\C\),
%given by \(\alpha\mapsto\langle \psi\mid A\alpha\rangle\), is linear
%and bounded. By the Riesz representation theorem, there is
%a unique \(\eta\in\Cal H\) such that
%\begin{equation}
%  f_{\psi}\circ A = f_{\eta} \equiv \langle\eta\mid-\rangle.
%\end{equation}
%Writing \(A^*\psi\equiv\eta\), this asignment gives us a
%bounded linear operator \(A^*\in\Cal L(\Cal H,\Cal H)\).
%Indeed, for every \(\psi\in\Cal H\) we have 
%\begin{align}
%  \|A^*\psi\|
%  &= \|\eta\| \\
%  &= \|f_\eta\| \\
%  &= \|f_\psi\circ A\| \\
%  &= \sup\{|\langle \psi\mid A\alpha\rangle| : \|\alpha\|\leq 1 \}\\
%  &\leq \sup\{\|\psi\|\|A\alpha\| : \|\alpha\|\leq 1 \}\\
%  &= \|\psi\|\|A\|.
%\end{align}
%This even proves that \(\|A^*\|\leq\|A\|\).
%In other words, if \(A\colon \Cal H\to\Cal H\) is \emph{bounded},
there is a unique bounded operator \(A^*\colon\Cal H\to\Cal H\),
called the \emph{adjoint} of \(A\), such that
\begin{equation}
  \forall \psi,\alpha\in\Cal H,\
  \langle \psi\mid A\alpha\rangle
  =
  \langle A^*\psi\mid \alpha\rangle.
\end{equation}
If moreover \(A=A^*\), we say that \(A\) is self-adjoint.

For various reasons, the observables in quantum mechanics are modeled
by \emph{self-adjoint operators}.
However, most of them are \emph{not bounded}.
It turns out that not-bounded operators defined on all of
\(\Cal H\) cannot be self-adjoint so, in order to have not-bounded
self-adjoint operators, we will need to allow their domain to be
only a \emph{dense subspace} \(\Cal D_A\subseteq\Cal H\).

\begin{definition}
  An \emph{operator on \(\Cal H\)} is a
  linear map $A\cl\mathcal{D}_A\to\mathcal{H}$,
  where $\mathcal{D}_A\subseteq\Cal H$ is a subspace.
  We say \(A\) is \emph{densely defined} if \(\Cal D_A\) is dense in
  \(\Cal H\), i.e.
  \begin{equation*}
    \forall \psi\in \mathcal{H}, \
    \forall \varepsilon > 0, \
    \exists \alpha \in \mathcal{D}_A, \
    \|\alpha - \psi\|< \varepsilon
  \end{equation*}
  or, equivalently, if for every $\psi\in\mathcal{H}$ there exists a
  sequence $\{\alpha_n\}_{n\in\N}$ in $\mathcal{D}_A$ whose limit is
  $\psi$.
\end{definition}
Now we define the \emph{adjoint} of a densely defined operator
\(A\colon\Cal D_A\to\Cal H\).
Note that we do \emph{not} assume that \(A\) is bounded, so when we
fix \(\psi\in\Cal H\), the linear functional
\begin{align*}
  f_\psi\circ A \colon \Cal D_A &\to \C \\
  \alpha &\mapsto \langle \psi \mid A\alpha \rangle
\end{align*}
is not necessarily bounded. However, if it \emph{is bounded},
by the BLT theorem, (\cref{thm:BLT}), it admits a unique
bounded extension \(\Cal H\to\C\) which, by the Riesz
representation theorem (\cref{thm:riesz}), is
of the form \(f_\eta\equiv\langle\eta\mid-\rangle\) for a unique
\(\eta\in\Cal H\).
In other words, \(\eta\) is the unique element of \(\Cal H\) such that
\begin{equation}
  f_\psi\circ A = f_{\eta}|_{\Cal D_A} \colon\Cal D_A\to\C
\end{equation}
or, equivalently,
\begin{equation}
  \forall\alpha\in\Cal D_A, \
  \langle \psi\mid A\alpha\rangle
  =
  \langle \eta\mid \alpha\rangle.
\end{equation}
Conversely, if there is some \(\eta\) satisfiyng these conditions,
then \(f_\psi\circ A\) is necessarily bounded (and \(\eta\) is unique).

Thus, we define the adjoint of \(A\) in the following way.
\begin{definition}
  Let $A\cl \mathcal{D}_A\to \mathcal{H}$ be a densely defined
  operator on $\mathcal{H}$. The \emph{adjoint}\index{adjoint} of $A$
  is the operator $A^*\cl\mathcal{D}_{A^*}\to\mathcal{H}$ defined as follows.
  \begin{enumerate}[label=(\roman*)]
    \item
      The domain of \(A^*\) is the subspace of \(\Cal H\) defined as
      \begin{align}
        \mathcal{D}_{A^*}
        &=
        \{
          \psi\in\mathcal{H}
          \mid
          f_\psi\circ A\colon\Cal D_A\to\C \text{ is bounded}
        \} \\
        &=
        \{
          \psi\in\mathcal{H}
          \mid
          \exists\eta\in\Cal H, \
          f_\psi\circ A=f_\eta|_{\Cal D_A}\colon\Cal D_A\to\C
        \} \\
        &=
        \{
          \psi\in\mathcal{H}
          \mid
          \exists\eta\in\Cal H, \
          \forall\alpha\in\Cal D_A, \
          \langle \psi\mid A\alpha \rangle = \langle \eta\mid\alpha\rangle
        \}
      \end{align}
    \item
      For each \(\psi\in\Cal D_{A^*}\) we define $A^*\psi$ to be the
      element \(\eta\) corresponding to \(\psi\).
      In other words, \(A^*\psi\) is
      the unique element of \(\Cal H\) satisfying
      \begin{equation}
        f_\psi\circ A = f_{A^*\psi}|_{\Cal D_A} \colon\Cal D_A\to\C.
      \end{equation}
      i.e.
      \begin{equation}
        \forall\alpha\in\Cal D_A, \
        \langle \psi\mid A\alpha\rangle
        =
        \langle A^*\psi\mid \alpha\rangle.
      \end{equation}
  \end{enumerate}
\end{definition}

\begin{remark}
  Although \(A\) is a densely defined operator,
  its adjoint \(A^*\) is not necessarily densely defined.
\end{remark}

If $A$ and $B$ are densely defined and $\mathcal{D}_A=\mathcal{D}_B$,
then the pointwise sum $A+B$ is clearly densely defined and hence,
$(A+B)^*$ exists. However, we do \emph{not} have $(A+B)^*=A^*+B^*$ in
general, unless one of $A$ and $B$ is bounded, but we do have the
following result.
\begin{proposition}
If $A$ is densely defined, then 
\begin{equation*}
(A+z\id_{\mathcal{D}_A})^*=A^*+\overline{z}\id_{\mathcal{D}_A}
\end{equation*}
for any $z\in\C$.
\end{proposition}
The identity operator $\id_{\mathcal{D}_A}$ is usually suppressed in the notation, so that the above equation reads $(A+z)^*=A^*+\overline{z}$.

% \begin{definition}
% Let $A\cl \mathcal{D}_A\to \mathcal{H}$ be a linear operator. The \emph{kernel}\index{kernel} and \emph{range}\index{range} of $A$ are
% \begin{equation*}
% \ker (A)  :=  \{\alpha \in \mathcal{D}_A \mid A\alpha = 0\},\quad\qquad\ran (A)  :=  \{A\alpha \mid \alpha \in \mathcal{D}_A \}.
% \end{equation*}
% The range is also called \emph{image}\index{image} and $\im (A)$ and $A(\mathcal{D}_A)$ are alternative notations.
% \end{definition}
% 
% \begin{proposition}
% An operator $A\cl \mathcal{D}_A\to \mathcal{H}$ is
% \begin{enumerate}[label=(\roman*)]
% \item injective if, and only if, $\ker(A)=\{0\}$
% \item surjective if, and only if, $\ran(A)=\mathcal{H}$.
% \end{enumerate}
% \end{proposition}
% 
% \begin{definition}
% An operator $A\cl \mathcal{D}_A\to \mathcal{H}$ is
% \emph{invertible}\index{invertible operator} if there exists an
% operator $B\cl\mathcal{H}\to \mathcal{D}_A$ such that $A\circ B =
% \id_{\mathcal{H}}$ and $B\circ A = \id_{\mathcal{D}_A}$.
% \end{definition}
% 
% \begin{proposition}
% An operator $A$ is invertible if, and only if, 
% \begin{equation*}
% \ker(A)=\{0\}\quad \text{ and }\quad\, \ran(A)=\mathcal{H}.
% \end{equation*}
% \end{proposition}

\begin{proposition}
\label{prp:kerran}
Let $A$ be densely defined. Then, $\ker(A^*)=\ran(A)^{\perp}$.
\end{proposition}

\begin{proof}
  Take \(\psi\in\ker(A^*)\), so \(A^*\psi=0\).
  Hence, for all \(\alpha\in\Cal D_A\), we have
  \begin{equation}
    \langle \psi \mid A\alpha \rangle
    = \langle A^*\psi \mid\alpha \rangle
    = 0
  \end{equation}
  which means that \(\psi\in\ran(A)^\perp\).
  Conversely, if \(\psi\in\ran(A)^\perp\), then
  \begin{equation}
    \langle \psi \mid A\alpha \rangle = 0.
  \end{equation}
  Hence, the linear functional \(f_\psi\circ A\colon\Cal D_A\to\C\) is
  bounded. Iin fact it is zero) and it is represen\(A^*\psi=0\), so
  \(\psi\in\ker(A^*)\).
\end{proof}

\begin{definition}
  A densely defined operator $A\cl\mathcal{D}_A\to\mathcal{H}$ is
  \emph{self-adjoint}\index{self-adjoint operator} if $A=A^*$. That is,
  \begin{enumerate}[label=(\roman*)]
    \item $\mathcal{D}_A = \mathcal{D}_{A^*}$
    \item $\forall \, \alpha \in \mathcal{D}_A : \ A\alpha = A^*\alpha$.
  \end{enumerate}
\end{definition}

\begin{definition}
  Let $A\cl \mathcal{D}_A \to \mathcal{H}$ and $B\cl \mathcal{D}_B \to
  \mathcal{H}$ be operators. We say that $B$ is an
  \emph{extension}\index{extension} of $A$, and we write $A\subseteq
  B$, iff
  \begin{enumerate}[label=(\roman*)]
  \item $\mathcal{D}_{A}\subseteq \mathcal{D}_{B}$
  \item $\forall \, \alpha \in \mathcal{D}_{A} : \ A\alpha = B \alpha$.
  \end{enumerate}
\end{definition}

\begin{proposition}
  \label{prp:adjointreverse}
  Let $A,B$ be densely defined. If $A\subseteq B$, then $B^*\subseteq
  A^*$.
\end{proposition}

\begin{proof}
  \begin{enumerate}[label=(\roman*)]
  \item Let $\psi \in \mathcal{D}_{B^*}$.
  Then the operator
  \begin{equation}
    f_\psi\circ B\colon \Cal D_B\to\C
  \end{equation}
  is bounded.
  Restricting to \(\Cal D_A\subseteq\Cal D_B\), we see that
  \begin{equation}
    f_\psi\circ A\colon \Cal D_A\to\C
  \end{equation}
  is bounded. Hence, \(\psi\in\Cal D_{A^*}\),
  which shows that $\mathcal{D}_{B^*}\subseteq \mathcal{D}_{A^*}$.
  \item
  If \(\psi\in\Cal D_{B^*}\subseteq\Cal D_{A^*}\), then for all
  \begin{equation}
    f_\psi\circ B = f_{B^*\psi}|_{\Cal D_B} \colon\Cal D_B\to\C.
  \end{equation}
  restricting to \(\Cal D_A\), we get
  \begin{equation}
    f_\psi\circ A = f_{B^*\psi}|_{\Cal D_A} \colon\Cal D_A\to\C,
  \end{equation}
  by uniqueness of \(A^*\psi\), we have $B^*\psi = A^*\psi$.
  \end{enumerate}
\end{proof}

\begin{corollary}
  \label[corollary]{cor:selfasdjext}
  A self-adjoint operator is maximal with respect to self-adjoint extension.
\end{corollary}

\begin{proof}
  Let $A,B$ be self-adjoint and suppose that $A\subseteq B$. Then
  \begin{equation*}
  A \subseteq B = B^* \subseteq A^* = A
  \end{equation*}
  and hence, $B=A$.
\end{proof}
%In fact, self-adjoint operators are maximal even with respect to
%symmetric extension, for we would have $B\subseteq B^*$ instead of
%$B=B^*$ in the above equation.

\section{The adjoint of a symmetric operator}

\begin{definition}
A densely defined operator $A\cl\mathcal{D}_A\to\mathcal{H}$ is called \emph{symmetric}\index{symmetric operator} if
\begin{equation*}
\forall \, \alpha,\beta \in \mathcal{D}_A : \ \langle\alpha|A\beta\rangle = \langle A\alpha|\beta\rangle.
\end{equation*}
\end{definition}

\begin{remark}
Observe that any self-adjoint operator is also symmetric, but a
symmetric operator need not be self-adjoint. 

In the physics literature, the term \emph{Hermitian}\index{Hermitian
operator} is sometimes used to refer to a symmetric operator, and
sometimes used to mean a self-adjoint operator.
To prevent confusions, we will avoid the term Hermitian altogether.
\end{remark}

%\begin{remark}
%In the physics literature, symmetric operators are usually referred to
%as \emph{Hermitian}\index{Hermitian operator} operators. However, this
%notion is then confused with that of self-adjointness when physicists
%say that observables in quantum mechanics correspond to Hermitian
%operators, which is not the case. On the other hand, if one decides to
%use Hermitian as a synonym of self-adjoint, it is then not true that
%all symmetric operators are Hermitian. In order to prevent confusion,
%we will avoid the term Hermitian altogether.
%\end{remark}

\begin{proposition}
  \label{prp:symext}
  A densely defined operator $A\colon\Cal D_A\to\Cal H$ is symmetric
  if, and only if, $A\subseteq A^*$.
\end{proposition}
\begin{proof}
  Suppose \(A\) is symmetric.
  Let $\psi\in\mathcal{D}_A$ and let $\eta:=A\psi$. Then, by symmetry, we have
  \begin{equation*}
    \forall \alpha\in \mathcal{D}_A, \
    \langle\psi|A\alpha\rangle = \langle \eta|\alpha\rangle
  \end{equation*}
  and hence $\psi\in\mathcal{D}_{A^*}$. Therefore,
  $\mathcal{D}_A\subseteq \mathcal{D}_{A^*}$ and $A^*\psi:=\eta =
  A\psi$.
  
  Conversely, suppose that \(A\subseteq A^*\). Then, for every
  \(\psi,\alpha\in\Cal D_A\) we have \(\psi\in D_{A^*}\) and
  \(A^*\psi=A\psi\), so
  \begin{equation}
    \langle \psi,A\alpha\rangle
    = \langle A^*\psi,\alpha\rangle
    = \langle A\psi,\alpha\rangle
  \end{equation}
  whence \(A\) is symmetric.
\end{proof}

\begin{corollary}
  A self-adjoint operator is maximal with respect to symmetric extensions.
\end{corollary}
\begin{proof}
  Same as \cref{cor:selfasdjext}, except we have \(B\subseteq B^*\)
  instead of \(B=B^*\).
\end{proof}

\section{Closability, closure, closedness}
\label{sec:closable-closure-closed}

\begin{definition}
  An operator \(A\colon\Cal D_A\to\Cal H\) (not necessarily densely defined) is \emph{closed} if its graph
  \begin{equation}
    \Gamma_A
    = \{(\psi,A\psi) : \psi\in\Cal D_A\}
    \subseteq \Cal H\times\Cal H
  \end{equation}
  is a closed subset of \(\Cal H\times\Cal H\).
  Equivalently, \(A\) is closed if, for each sequence \(\psi_n\) in
  \(\Cal D_A\) with \(\psi_n\to\psi\in\Cal H\) and \(A\psi_n \to
  \alpha\in\Cal H\), we have \(\psi\in\Cal D_A\) and \(A\psi=\alpha\).

  The operator \(A\) is called \emph{closable} if the closure of its
  graph is the graph of an operator \(\overline A\):
  \begin{equation}
    \overline{\Gamma_A}
    = \{(\psi,\overline A\psi) : \psi\in\Cal D_{\overline A}\}
    \subseteq \Cal H\times\Cal H.
  \end{equation}
  Then \(\overline A\) is necessarily unique, it is closed and it
  extends \(A\).
\end{definition}

\begin{remark}
Note that we have used the overline notation in several contexts with
different meanings. When applied to complex numbers, it denotes
complex conjugation. When applied to subsets of a topological space,
it denotes their topological closure. Finally, when applied to
(closable) operators, it denotes their closure as defined above.
\end{remark}

\begin{proposition}
  If \(A\colon\Cal D_A\to\Cal H\) is a densely defined operator, then
  its adjoint \(A^*\) is closed.
\end{proposition}
\begin{proof}
  Suppose \(\psi_n\in\Cal D_{A^*}\) converges to \(\psi\in\Cal H\) and
  \(A^*\psi_n\) converge to \(\alpha\in\Cal H\).
  Then, for every \(\beta\in\Cal D_A\), we have
  \begin{align}
    \langle \alpha \mid \beta \rangle
    &= \lim_{n\to\infty}\langle A^*\psi_n \mid \beta \rangle \\
    &= \lim_{n\to\infty}\langle \psi_n \mid A\beta \rangle \\
    &= \langle \psi \mid A\beta \rangle.
  \end{align}
  Hence, the operators \(f_\psi\circ A\colon \Cal D_A\to\C\) and
  \(f_\alpha|_{\Cal D_A}\colon\Cal D_A\to\C\) are equal,
  which means that \(\psi\in\Cal D_{A^*}\) and \(A^*\psi=\alpha\).
\end{proof}

\begin{proposition}
  \label{lem:closableext}
  Let \(A\colon\Cal D_A\to\Cal H\) be a densely defined operator.
  \begin{enumerate}[label=(\roman*)]
    \item 
      If $A$ is closable, then \(A^*\) is densely defined.
    \item
      Conversely, if $A^*$ is densely defined,
      then \(A\subseteq A^{**}\), so \(A\) is closable and, in fact,
      $\overline{A}=A^{**}$.
  \end{enumerate}
\end{proposition}

\begin{proof}
  First we check (i).
  Suppose that \(A^*\) is densely defined.
  We begin by showing that \(\Cal D_A\) is contained in
  \begin{equation}
    \mathcal{D}_{A^{**}}=
    \{\psi\in \mathcal{H}\mid
    \exists\eta\in\Cal H,\
    \forall\alpha\in \mathcal{D}_{A^*}, \
    \langle \psi | A^*\alpha \rangle=\langle \eta | \alpha \rangle
    \}
  \end{equation}
  Take \(\psi\in \Cal D_A\). By definition of \(A^*\),
  \begin{equation*}
    \forall \alpha \in\mathcal{D}_{A^*}, \
    \langle \alpha | A\psi \rangle =  \langle A^* \alpha | \psi \rangle.
  \end{equation*}
  Taking complex conjugates, we obtain
  \begin{equation*}
    \forall \alpha \in\mathcal{D}_{A^*}, \
    \langle \psi | A^* \alpha \rangle
    = 
    \langle A\psi | \alpha \rangle,
  \end{equation*}
  so \(\psi\in\Cal D_{A^{**}}\).
  This also shows that $A^{**}\psi= A\psi$, so $A\subseteq A^{**}$, as
  we wanted. To show that \(A^{**}=\overline{A}\),
  take anoter closed extension \(B\) of \(A\).
  Since \(A\subseteq B\), then \(B\) is necessarily densely defined,
  so applying \cref{prp:adjointreverse} we obtain
  \(B^*\subseteq A^*\).
  twice, we see that
  \(A^{**}\subseteq B^{**}=B\), as we wanted.
  \end{proof}
  
  \begin{proposition}
  Let \(A\) be a symmetric operator. Then
  \begin{enumerate}[label=(\roman*)]
    \item \(A\) is closable,
    \item \(A^{**}\subseteq A^*\), so \(A\subseteq A^{**}\subseteq A^*\),
    \item \(A^{**}\) is symmetric.
  \end{enumerate}
  \end{proposition}
\begin{proof}
  By definition, a symmetric operator \(A\) is densely defined.
  Moreover, in \cref{prp:symext}, we saw that $A\subseteq A^*$.
  In particular, $\mathcal{D}_A\subseteq \mathcal{D}_{A^*}$, so
  \begin{equation*}
    \mathcal{H}
    =\overline{\mathcal{D}_A}\subseteq \overline{\mathcal{D}_{A^*}}.
  \end{equation*}
  Hence, $\overline{\mathcal{D}_{A^*}} = \mathcal{H}$, so \(A\) is
  closable. 
  For (ii), note that, since $A$ is symmetric, we have
  $A\subseteq A^*$ by \cref{prp:symext}.
  Hence, $A^{**}\subseteq A^*$ by \cref{prp:adjointreverse}.
  Finally, for (iii), take \(\alpha,\psi\in\Cal D_{A^{**}}\).
  By part (ii), we have \(\psi\in\Cal D\) and \(A^{**}\psi=A^*\psi\),
  so
  \begin{equation}
    \langle A^{**}\alpha\mid\psi\rangle
    = \overline{\langle A^*\psi\mid\alpha\rangle}
    = \langle \alpha\mid A^{**}\psi\rangle
  \end{equation}
\end{proof}

Note that the adjoint of a symmetric operator need
\emph{not} be symmetric.

% What?
\begin{theorem}
Let $A$ be a densely defined operator.
\begin{enumerate}[label=(\roman*)]
\item If $A$ is invertible, we have $(A^{-1})^*=(A^*)^{-1}$
\item If $A$ is invertible and closable and $\overline{A}$ is injective, then $\overline{A^{-1}}=\overline{A}^{\,-1}$.
\end{enumerate}
\end{theorem}

\section{Essentially self-adjoint operators}

Usually, checking that an operator is symmetric is easy. By contrast, checking that an operator is self-adjoint (directly from the definition) requires the construction of the adjoint, which is not always easy. However, since every self-adjoint operator is symmetric, we can first check the symmetry property, and then determine criteria for when a symmetric operator is self-adjoint, or for when a self-adjoint extension exists.

Two complications with the extension approach are that, given a symmetric operator, there could be \emph{no} self-adjoint extension at all, or there may be \emph{several} different self-adjoint extensions. There is, however, a class of operators for which the situation is much nicer.

\begin{definition}
A symmetric operator $A$ is called \emph{essentially self-adjoint}\index{essentially self-adjoint operator} if $\overline{A}$ is self-adjoint.
\end{definition}

This is weaker than the self-adjointness condition.

\begin{proposition}
If $A$ is self-adjoint, then it is essentially self-adjoint. 
\end{proposition}

\begin{proof}
If $A = A^*$, then $A\subseteq A^*$ and $A^* \subseteq A$. Hence, $A^{**}\subseteq A^*$ and $A^*\subseteq A^{**}$, so $A^* = A^{**}$. Similarly, we have $A^{**} = A^{***}$, which is just $\overline{A}=\overline{A}^{\,*}$.
\end{proof}

\begin{theorem}
If $A$ is essentially self-adjoint, then there exists a unique self-adjoint extension of $A$, namely $\overline{A}$.
\end{theorem}

\begin{proof}
\begin{enumerate}[label=(\roman*)]
\item Since $A$ is symmetric, it is closable and hence $\overline{A}$ exists.
\item By \Cref{lem:closableext}, we have $A\subseteq \overline{A}$, so $\overline{A}$ is an extension of $A$. 
\item Suppose that $B$ is another self-adjoint extension of $A$. Then, $A\subseteq B = B^*$, hence $B^{**}\subseteq A^*$, and thus $A^{**}\subseteq B^{***}=B$. This means that $\overline{A}\subseteq B$, i.e.\ $B$ is a self-adjoint extension of the self-adjoint operator $\overline{A}$. Hence, $B=\overline{A}$ by \Cref{cor:selfasdjext}.\qedhere
\end{enumerate}
\end{proof}

\begin{remark}
One may get the feeling at this point that checking for essential self-adjointness of an operator $A$, i.e.\ checking that $A^{**}=A^{***}$, is hardly easier than checking whether $A$ is self-adjoint, that is, whether $A=A^*$. However, this is not so. While we will show below that there is a sufficient criterion for self-adjointness which does not require to calculate the adjoint, we will see that there is, in fact, a necessary and sufficient criterion to check for essential self-adjointness of an operator without calculating a single adjoint.
\end{remark}

\begin{remark}
If a symmetric operator $A$ fails to even be essentially self-adjoint, then there is either no self-adjoint extension of $A$ or there are several. 
\end{remark}

\begin{definition}
Let $A$ be a densely defined operator. The \emph{defect indices}\index{defect indices} of $A$ are
\begin{equation*}
d_+:=\dim (\ker(A^*-\mathrm{i})) , \qquad \quad d_- :=\dim (\ker (A^*+\mathrm{i})),
\end{equation*}
where by $A^*\pm\mathrm{i}$ we mean, of course, $A^*\pm\mathrm{i}\cdot\id_{\mathcal{D}_A}$.
\end{definition}

\begin{theorem}
A symmetric operator has a self-adjoint extension if its defect indices coincide. Otherwise, there exist no self-adjoint extension.
\end{theorem}

\begin{remark}
We will later see that if $d_+=d_-=0$, then $A$ is essentially self-adjoint.
\end{remark}

\section{\texorpdfstring
  {Criteria for self-adjointness \\ and essential self-adjointness}
  {Criteria for self-adjointness and essential self-adjointness}
}

\begin{theorem}
A symmetric operator $A$ is self-adjoint if (but not only if)
\begin{equation*}
\exists \, z\in\C :\ \ran(A+z)=\mathcal{H}=\ran(A+\overline{z}).
\end{equation*}
\end{theorem}

\begin{proof}
Since $A$ is symmetric, by \Cref{prp:symext}, we have $A\subseteq A^*$. Hence, it remains to be shown that $A^*\subseteq A$. To that end, let $\psi\in \mathcal{D}_{A^*}$ and let $z\in \C$. Clearly, 
\begin{equation*}
A^*\psi+\overline{z}\psi\in\mathcal{H}.
\end{equation*}
Now suppose that $z$ satisfies the hypothesis of the theorem. Then, as $\ran(A+\overline{z})=\mathcal{H}$,
\begin{equation*}
\exists \, \alpha\in\mathcal{D}_{A} : \ A^*\psi+\overline{z}\psi = (A+\overline{z})\alpha.
\end{equation*}
By using the symmetry of $A$, we have that, for any $\beta\in\mathcal{D}_A$,
\begin{IEEEeqnarray*}{rCl}
\langle \psi| (A+z)\beta \rangle & = & \langle (A+z)^*\psi| \beta \rangle \\
 & = & \langle A^*\psi+\overline{z}\psi| \beta \rangle \\
 & = & \langle (A+\overline{z})\alpha | \beta \rangle  \\
 & = & \langle A\alpha | \beta \rangle +z\langle \alpha | \beta \rangle \\
 & = & \langle \alpha | A\beta \rangle +\langle \alpha | z\beta \rangle \\
 & = & \langle \alpha | (A+z)\beta \rangle.
\end{IEEEeqnarray*}
Since $\beta\in\mathcal{D}_A$ was arbitrary and $\ran(A+z)=\mathcal{H}$, we have
\begin{equation*}
\forall \, \varphi \in \mathcal{H} : \quad \langle \psi | \varphi\rangle= \langle \alpha | \varphi \rangle.
\end{equation*}
Hence, by positive-definiteness of the inner product, we have $\psi\in\alpha$, thus $\psi\in\mathcal{D}_A$ and therefore, $A^*\subseteq A$.
\end{proof}


\begin{theorem}
A symmetric operator $A$ is essentially self-adjoint if, and only if, 
\begin{equation*}
\exists \, z\in \C\setminus \R : \ \overline{\ran(A+z)}=\mathcal{H}=\overline{\ran(A+\overline{z})}. 
\end{equation*}
\end{theorem}

The following criterion for essential self-adjointness, which does require the calculation of $A^*$, is equivalent to the previous result and, in some situations, it can be easier to check.

\begin{theorem}
A symmetric operator $A$ is essentially self-adjoint if, and only if,
\begin{equation*}
\exists \, z\in \C\setminus \R : \ \ker(A^*+z)=\{0\}=\ker(A^*+\overline{z}). 
\end{equation*}
\end{theorem}

\begin{proof}
We show that this is equivalent to the previous condition. Recall that if $\mathcal{M}$ is a linear subspace of $\mathcal{H}$, then $\mathcal{M}^{\perp}$ is closed and hence, $\mathcal{M}^{\perp\perp}=\overline{M}$. Thus, by \Cref{prp:kerran}, we have
\begin{equation*}
\overline{\ran(A+z)}  =  \ran(A+z)^{\perp\perp}
 =  \ker(A^*+\overline{z})^{\perp}
\end{equation*}
and, similarly,
\begin{equation*}
\overline{\ran(A+z)} = \ker(A^*+z).
\end{equation*}
Since $\mathcal{H}^{\perp}=\{0\}$, the above condition is equivalent to
\begin{equation*}
\overline{\ran(A+z)}=\mathcal{H}=\overline{\ran(A+\overline{z})}.\qedhere
\end{equation*}
\end{proof}





















